% !TEX root = ../matan.tex

\section{Признак Коши}

Имеется в виду \textbf{радикальный признак Коши} (корневой признак).

\begin{Theorem}{Признак Коши}{}
	Пусть $\sum_{k=1}^{\infty} a_k$ --- произвольный ряд, $|a_k| > 0$ для
	всех $k$. Если существуют $k_0 \in \mathbb{N}$ и $q < 1$ такие, что
	\[
	(|a_k|)^{1/k} \le q
	\quad \text{для всех } k \ge k_0,
	\]
	то ряд $\sum_{k=1}^{\infty} a_k$ сходится.
\end{Theorem}

\begin{proof}
	Из условия $(|a_k|)^{1/k} \le q$ при $k \ge k_0$ следует $|a_k| \le q^k$.
	Запишем частичную сумму:
	\[
	S_n = \sum_{k=1}^{n} a_k
	= \underbrace{\sum_{k=1}^{k_0-1} a_k}_{\text{фикс. конечная сумма}}
	+ \sum_{k=k_0}^{n} a_k.
	\]
	Первая часть --- фиксированное конечное число. Оценим хвост:
	\[
	\sum_{k=k_0}^{n} |a_k|
	\;\le\; \sum_{k=k_0}^{n} q^k
	\;\le\; \frac{q^{k_0}}{1-q}
	\;<\; +\infty,
	\]
	так как $q < 1$ и $\sum q^k$ --- сходящийся геометрический ряд. Частичные
	суммы $\sum_{k=k_0}^{n} |a_k|$ монотонны и ограничены, значит ряд
	$\sum_{k=k_0}^{\infty} |a_k|$ сходится. Абсолютная сходимость влечёт
	сходимость, поэтому $\sum_{k=k_0}^{\infty} a_k$ сходится. Вместе с конечной
	начальной частью ряд $\sum_{k=1}^{\infty} a_k$ сходится.
\end{proof}

\begin{Corollary}{Признак Коши в форме верхнего предела}{}
	Если $\displaystyle\limsup_{k\to\infty}(|a_k|)^{1/k} \le q < 1$, то ряд
	$\sum_{k=1}^{\infty} a_k$ сходится.
\end{Corollary}

\begin{proof}
	Пусть $\limsup_{k\to\infty}(|a_k|)^{1/k} \le q < 1$. Выберем $r$ с $q < r < 1$.
	Так как верхний предел $\le q < r$, по определению верхнего предела найдётся
	$k_0$ такое, что $(|a_k|)^{1/k} \le r$ для всех $k \ge k_0$. Это в точности
	условие основной теоремы с $r$ вместо $q$, откуда ряд $\sum a_k$ сходится.
\end{proof}

\begin{Remark}{Случай $q = 1$ неопределён}{}
	Если $(|a_k|)^{1/k} \to 1$, то признак ответа не даёт: для $a_k = \dfrac{1}{k}$
	имеем $\sqrt[k]{1/k} \to 1$, но $\sum\dfrac{1}{k}$ расходится; для
	$a_k = \dfrac{1}{k^2}$ также $\sqrt[k]{1/k^2} \to 1$, но $\sum\dfrac{1}{k^2}$
	сходится. (Здесь использовано $\sqrt[k]{k} \to 1$ при $k \to \infty$.)
\end{Remark}

\begin{Remark}{Сравнение с признаком Даламбера}{}
	Радикальный признак Коши не слабее признака Даламбера. Более точно, всегда
	\[
	\varliminf_{k\to\infty}\frac{a_{k+1}}{a_k}
	\;\le\; \varliminf_{k\to\infty}\sqrt[k]{a_k}
	\;\le\; \varlimsup_{k\to\infty}\sqrt[k]{a_k}
	\;\le\; \varlimsup_{k\to\infty}\frac{a_{k+1}}{a_k}.
	\]
	Поэтому если предел отношения $\dfrac{a_{k+1}}{a_k}$ существует и равен $q$,
	то и корневой предел существует и равен $q$. Но существуют ряды, к которым
	применим признак Коши, а признак Даламбера --- нет.
\end{Remark}